% ---------------------------------------------------------------------------
% Conformance case for the PUBLIC API (Protocol A).
%
% The point of this case is not that these particular examples render: it
% is that a syntax front-end can be built WITHOUT the package's internals.
% Everything below the \ExplSyntaxOn uses \lx_... names only -- no \lx@...,
% no \makeatletter, no assumption about how the machinery works.  If a
% later change breaks the API, this case stops compiling; if a later change
% breaks the SEMANTICS the API promises, the assertions in runtests.py
% (a_frontend) catch it.
%
% Two front-end shapes are exercised, because they use the API differently:
%
%   \pex ... blank line          the self-contained shape: the example
%                                brings its own list.  This is how \ex.
%                                is built.
%   \begin{pexe} \pit ... \end   the batch shape: one list holds several
%                                examples.  This is how exe is built.
%
% The preamble is not shared with the other cases: _preamble-ua.tex already
% contains \begin{document}, and the front-end has to be defined before
% that.  It mirrors _preamble-ua otherwise, and for the same reason -- the
% API hands out the list funnel and the tagged-Span helpers, so the promise
% worth checking is not merely that the page looks right but that what a
% front-end builds out of them is VALID PDF/UA.  veraPDF is the only
% authoritative oracle for that, and a_frontend runs it.
%
% "uncompress" is added on top of the UA settings so the structure-tree
% helpers can read the tree out of the raw bytes; it does not affect
% validity.  A dc:title is mandatory under ISO 14289-2 clause 8.11.1,
% hence the pdftitle.  Three LaTeX passes are needed for UA convergence
% under pdflatex and xelatex -- see PASSES in runtests.py.
% ---------------------------------------------------------------------------
\DocumentMetadata{uncompress,lang=en,pdfversion=2.0,pdfstandard=ua-2,tagging=on}
\documentclass{article}
\usepackage{iftex}
\ifPDFTeX
  \usepackage[T1]{fontenc}
  \usepackage[utf8]{inputenc}
\else
  \usepackage{fontspec}
\fi
\usepackage{linguexx}
\usepackage{hyperref}
\hypersetup{pdftitle={linguexx front-end API conformance case},
            pdfdisplaydoctitle=true}
\pagestyle{empty}

\ExplSyntaxOn

%% ---- shape 1: the self-contained example -------------------------------
%% \pex collects its body under the package's own terminator rules (blank
%% line, \z., unmatched \end) and hands it to a continuation, which is
%% nothing but: open the example, make an item, typeset, close.
\NewDocumentCommand \pex { } { \lx_body_collect:N \__fe_run:n }
\cs_new_protected:Npn \__fe_run:n #1
  {
    \lx_example_begin:
    \lx_item_main:
    #1
    \lx_example_end:
  }

%% the same with a label of the front-end's choosing (no counter step)
\tl_new:N \l__fe_label_tl
\NewDocumentCommand \pexc { m }
  {
    \tl_set:Nn \l__fe_label_tl {#1}
    \lx_body_collect:N \__fe_runc:n
  }
\cs_new_protected:Npn \__fe_runc:n #1
  {
    \lx_example_begin:
    \exp_args:NV \lx_item_main:n \l__fe_label_tl
    #1
    \lx_example_end:
  }

%% ...and one whose judgment is an ARGUMENT rather than something typed
%% into the text -- the case that cannot be composed out of the scanning
%% commands, since \lx_judgment_scan: clears the mark before it peeks.
%%
%% Note the ORDER, which is the API's first invariant: \lx_item_core: sets
%% the label, and only then is the list opened, because under [legacy] the
%% list geometry is read off the label.  Opening the list first renders
%% identically in the default mode and mis-sizes the label box in legacy,
%% which is exactly the kind of bug this case exists to prevent.
\NewDocumentCommand \pexj { m }
  {
    \tl_set:Nn \l__fe_label_tl {#1}
    \lx_body_collect:N \__fe_runj:n
  }
\cs_new_protected:Npn \__fe_runj:n #1
  {
    \lx_example_begin:
    \lx_item_core:
    \lx_list_main_open:
    \exp_args:NV \lx_judgment_set:n \l__fe_label_tl
    \lx_item_emit:
    #1
    \lx_example_end:
  }

%% sub-levels: push opens a new, deeper one; next stays at the current one
\NewDocumentCommand \pa { } { \lx_sub_push: }
\NewDocumentCommand \pb { } { \lx_sub_next: }

%% ---- shape 2: the batch ------------------------------------------------
%% One list, several examples.  The list is opened before any item exists,
%% so the label it is sized from has to be guessed.
\NewDocumentEnvironment { pexe } { }
  {
    \lx_example_begin:
    \lx_label_guess:
    \lx_list_main_open:
  }
  { \lx_example_end: }
\NewDocumentCommand \pit { } { \lx_item_next: }
%% In the batch the list is already open, so the bundled form applies:
%% label, mark and item in one move, at whatever level is current.
\NewDocumentCommand \pitj { m } { \lx_item_judged:n {#1} }

\ExplSyntaxOff

\begin{document}
\section{FESECT a front-end built on the public API}

FEINTRO ordinary running text precedes the first example.

\pex FEPLAIN one.

\pex *FEJUDGED two.

\pex FESUBHOST three.
\pa FESUBA alpha.
\pb FESUBB beta.
\pa FEROMANA roman.

% A SECOND example with sub-levels: its letters must restart at a., which
% is the only thing that shows \lx_example_begin: resetting the sub-counters.
\pex FESUBHOSTB four.
\pa FESUBC gamma.
\pb FESUBD delta.

\pexc{(fex)} FECUSTOM five.

\pexj{\#} FEARGJUDGE six.

\begin{pexe}
\pit FEBATCHONE seven.
\pit FEBATCHTWO eight.
\pitj{??} FEBATCHJ nine.
\end{pexe}

\pex FEAFTER ten.

FEPARA A paragraph after the examples.

\noindent FEMARGIN At the margin.

\end{document}
